\clearpage
\section{Selecting primary and secondary chains}
\label{sec:chaining-primary-secondary}

After chains are collected, we must note that several chains can naturally match the same part of a query in differents parts of the target.
Sometimes there can be significant overlap between better chain and slightly worse one. For this reason minimap2 needs to classify chains 
into two groups: primary chains and secondary chains and it reports all primary chains but limit reported secondary chains to some set runtime limit~\autocite{mmseqs2_issue_323}.
Minimap considers chains from the highest score to the lowest:
\begin{itemize}
  \item Chains are classified either as \emph{primary} (best one) or \emph{secondary}
  \item Specific chain becomes \emph{secondary} when there's existing \emph{primary} that covers enough of the same query interval 
  \item Otherise, chain is classified as \emph{primary}
\end{itemize}
BioseqDB uses the same classification logic as minimap2. Below, in Algorithm~\ref{alg:chaining-primary-secondary}, I show how this classification is formally done:
\begin{itemize}
  \item $c$ is a chain (iterated by decreasing score)
  \item \(I_c=[q_c^{\mathrm{start}},q_c^{\mathrm{end}})\) is the query interval of chain $c$ (start of first anchor and end of last anchor).
  We will intersect these intervals to calcualte how many positions chain covers in the query.
  \item \(s_c\) is chain score
  \item The value \(u\) counts symbols (bases) not covered by \emph{any} accepted primary chain.
  \item Previously $p$ was specifically parent of an anchor. In this case it represents primary chain.
  \item $\pi(c)$ is the parent chain of $c$ (I defined previously parents for anchors, but this is extended notation to represent parent of first anchor in chain)
  \item Subtracting this contribution makes a chain with substantial new coverage less likely to become secondary.
  \item The default overlap threshold is \(\lambda=0.5\).
  \item The uncovered-length limit is \(L=2^{31}-1\), so it normally adds no restriction.
\end{itemize}

Let \(I_c=[q_c^{\mathrm{start}},q_c^{\mathrm{end}})\) be the query interval
of chain \(c\), and \(s_c\) its score. Below, in
Algorithm~\ref{alg:chaining-primary-secondary}, I show how classification is formally done:\(|I|\) counts query bases.
The value \(u\) counts bases not covered by \emph{any} accepted primary.
Subtracting this contribution makes a chain with substantial new coverage
less likely to become secondary. The default overlap threshold is
\(\lambda=0.5\). The uncovered-length limit is \(L=2^{31}-1\), so it
normally adds no restriction.

\begin{algorithm}[H]
\caption{Minimap2 primary and secondary chain selection}
\label{alg:chaining-primary-secondary}
\small
\begin{algorithmic}[1]
  \Require chains \(C\), query intervals \(I_c\), scores \(s_c\), limits \(\lambda,L\)
  \Ensure ordered primary list \(P\) and parent \(\pi(c)\) for every chain
  \State \(P\gets[\,]\)
  \For{\(c\in C\) in decreasing order of \(s_c\)}
    \State \(u\gets\bigl|I_c\setminus\bigcup_{p\in P}I_p\bigr|,\quad \pi(c)\gets c\)
      \Comment{initially primary}
    \For{\(p\in P\) in acceptance order}
      \State \(o\gets|I_c\cap I_p|\)
      \State \(a\gets\min(|I_c|,|I_p|),\quad b\gets\max(|I_c|,|I_p|)\)
      \If{\(o>0\land o/a-u/b>\lambda\land u\leq L\)}
        \State \(\pi(c)\gets p\) \Comment{secondary alternative to \(p\)}
        \State \textbf{break}
      \EndIf
    \EndFor
    \If{\(\pi(c)=c\)}
      \State append \(c\) to \(P\)
    \EndIf
  \EndFor
\end{algorithmic}
\end{algorithm}

A separate filter retains all primaries and limits secondary alternatives.
With the defaults, a secondary must score at least \(0.8s_{\pi(c)}\) or be
within \(2k\) of its parent's score, where \(k\) is the seed length.
Exact duplicates of the parent's coordinates are removed, and at most five
secondaries survive across the list. BioseqDB follows these selection steps
and repeats them after base alignment refines the regions
(Section~\ref{sec:alignment-postprocessing}).

Primary \emph{chains} should not be confused with primary SAM records.
Minimap2 marks one surviving nonsecondary alignment as the primary SAM record.
The other nonsecondary parts of a split mapping are \emph{supplementary}.
Secondary records instead describe alternative mappings of the same query
region.
\clearpage
